% !TeX program = pdflatex
% =============================================================================
% COPYRIGHT AND DISCLAIMER
% =============================================================================
% This file is provided as a convenience template for competition participants.
% By using it, you accept full responsibility for compliance with all applicable
% copyright, licensing, and usage rules of the Association for Computing
% Machinery (ACM) and of any LaTeX packages you load (including, without
% limitation, the "acmart" document class and related ACM materials).
%
% Consult ACM's official publication, template, and copyright policies before
% using this setup for any purpose other than this competition. Do not assume
% that this template may be reused, redistributed, or submitted elsewhere
% without independently verifying ACM's current rules.
%
% Use of the "acmart" class, the "sigconf" option, or ACM-style formatting in
% this template does NOT mean that this report, this competition, its organizers,
% or any resulting PDF is sponsored, endorsed, affiliated with, or published by
% ACM. ACM trademarks and branding remain the property of ACM.
%
% Template packages live in preamble.tex.
% Citations live in main.bib. Edit this file with your paper content.
% =============================================================================

\documentclass[sigconf]{acmart}
% =============================================================================
% Template setup — participants normally do not edit this file.
% Author content belongs in main.tex; citations belong in main.bib.
% \documentclass lives in main.tex so Overleaf does not treat this file as root.
% =============================================================================

% --- Competition / non-ACM-submission friendly rights metadata -------------
\setcopyright{none}
\settopmatter{printacmref=false}
\acmConference[Kaggle Competition]{Filament Segmentation Challenge}{2026}{\href{https://www.kaggle.com/competitions/filament-segmentation-2026}{kaggle.com/competitions/filament-segmentation-2026}}
\acmBooktitle{Competition Report}
\acmDOI{}
\acmISBN{}

% --- Author guidance text --------------------------------------------------
% \guide{...} is gray italic while \guideStyletrue; omitted from the PDF if false.
% Flip the switch in main.tex after % =============================================================================
% Template setup — participants normally do not edit this file.
% Author content belongs in main.tex; citations belong in main.bib.
% \documentclass lives in main.tex so Overleaf does not treat this file as root.
% =============================================================================

% --- Competition / non-ACM-submission friendly rights metadata -------------
\setcopyright{none}
\settopmatter{printacmref=false}
\acmConference[Kaggle Competition]{Filament Segmentation Challenge}{2026}{\href{https://www.kaggle.com/competitions/filament-segmentation-2026}{kaggle.com/competitions/filament-segmentation-2026}}
\acmBooktitle{Competition Report}
\acmDOI{}
\acmISBN{}

% --- Author guidance text --------------------------------------------------
% \guide{...} is gray italic while \guideStyletrue; omitted from the PDF if false.
% Flip the switch in main.tex after % =============================================================================
% Template setup — participants normally do not edit this file.
% Author content belongs in main.tex; citations belong in main.bib.
% \documentclass lives in main.tex so Overleaf does not treat this file as root.
% =============================================================================

% --- Competition / non-ACM-submission friendly rights metadata -------------
\setcopyright{none}
\settopmatter{printacmref=false}
\acmConference[Kaggle Competition]{Filament Segmentation Challenge}{2026}{\href{https://www.kaggle.com/competitions/filament-segmentation-2026}{kaggle.com/competitions/filament-segmentation-2026}}
\acmBooktitle{Competition Report}
\acmDOI{}
\acmISBN{}

% --- Author guidance text --------------------------------------------------
% \guide{...} is gray italic while \guideStyletrue; omitted from the PDF if false.
% Flip the switch in main.tex after \input{preamble}.
%
% Do not use \guide{...} for \email{...}.
\usepackage{xcolor}

% --- Math ------------------------------------------------------------------
% Do NOT load amssymb (or amsfonts): acmart already loads newtxmath, which
% defines the same symbols (e.g. \Bbbk) and causes a conflict on Overleaf.
% amsmath is already provided by acmart; mathtools extends it.
\usepackage{mathtools}

% --- Figures ---------------------------------------------------------------
% graphicx is already loaded by acmart
\usepackage{subcaption}

% --- Tables ----------------------------------------------------------------
% booktabs is already loaded by acmart
\usepackage{tabularx}
\usepackage{multirow}

% --- Bibliography style is set in main.tex (ACM-Reference-Format) ----------
.
%
% Do not use \guide{...} for \email{...}.
\usepackage{xcolor}

% --- Math ------------------------------------------------------------------
% Do NOT load amssymb (or amsfonts): acmart already loads newtxmath, which
% defines the same symbols (e.g. \Bbbk) and causes a conflict on Overleaf.
% amsmath is already provided by acmart; mathtools extends it.
\usepackage{mathtools}

% --- Figures ---------------------------------------------------------------
% graphicx is already loaded by acmart
\usepackage{subcaption}

% --- Tables ----------------------------------------------------------------
% booktabs is already loaded by acmart
\usepackage{tabularx}
\usepackage{multirow}

% --- Bibliography style is set in main.tex (ACM-Reference-Format) ----------
.
%
% Do not use \guide{...} for \email{...}.
\usepackage{xcolor}

% --- Math ------------------------------------------------------------------
% Do NOT load amssymb (or amsfonts): acmart already loads newtxmath, which
% defines the same symbols (e.g. \Bbbk) and causes a conflict on Overleaf.
% amsmath is already provided by acmart; mathtools extends it.
\usepackage{mathtools}

% --- Figures ---------------------------------------------------------------
% graphicx is already loaded by acmart
\usepackage{subcaption}

% --- Tables ----------------------------------------------------------------
% booktabs is already loaded by acmart
\usepackage{tabularx}
\usepackage{multirow}

% --- Bibliography style is set in main.tex (ACM-Reference-Format) ----------


\guideStylefalse

\begin{document}

\title{Beyond Annotator Agreement: A Calibrated Three-Family Ensemble with Crop Refinement for Solar Filaments}
\subtitle{A Solution to the Solar Filament Segmentation Challenge 2026}


\author{Vedran Hrabar}
\affiliation{
  \institution{Faculty of Electrical Engineering and Computing, University of Zagreb}
  \city{Zagreb}
  \country{Croatia}
}
\email{vedran.hrabar@fer.unizg.hr}

\begin{abstract}
    This report describes solution to the Solar Filament Segmentation Challenge~2026, a Kaggle competition on automatic segmentation of solar filaments in GONG H-$\alpha$ observations.
\end{abstract}

\maketitle


% ==================================== %
%              INTRODUCTION            %
% ==================================== %
\section{Introduction}\label{sec:introduction}
    This report goes over the details of the submitted solution to the Solar Filament Segmentation Challenge~2026~\cite{filament-segmentation-2026}.
    The challenge, launched on July~10,~2026, is organized by the Earth-Space AI Research (ESAIR) Lab\footnote{\href{https://www.esairlab.com/home}{www.esairlab.com/}}.
    This challenge utilizes the MAGFiLO dataset for evaluation of the segmentation algorithms~\cite{ahmadzadeh2024dataset}.


% ==================================== %
%      DATA AND VALIDATION PROTOCOL    %
% ==================================== %
\section{Data and Validation Protocol}\label{sec:data}
    \textbf{Labels.} The training set contains 707 observations and 1154 annotation records (411 observations with one annotator, 145 with two, 151 with three).
    Together, these records hold 8199 filament polygons with a median area of 1228 px.
    Thin structures under about 5 px, i.e.\ barbs, account for only 1.9\% of the annotated area.
    One annotator's set was scored against another's on the 296 observations with several records.
    This comparison gives PQ~0.34 (SQ 0.634, RQ 0.54), and annotators agree on barb pixels only 40\% of the time.
    Most disagreement concerns whether a faint filament exists and where a long one ends.
    A score of 0.47 is reached by the final model against the same records, so it agrees with each annotator more than the annotators agree with each other.

    \textbf{Scoring.} PQ~\cite{kirillov2019panoptic} matches a prediction to a filament at IoU~$>0.5$ and computes $\mathrm{PQ}=\sum_{\mathrm{TP}}\mathrm{IoU}\,/\,(|\mathrm{TP}|+\tfrac12|\mathrm{FP}|+\tfrac12|\mathrm{FN}|)$.
    Ground-truth instances are the individual polygons rasterised as COCO RLE~\cite{lin2014coco}, and each record is scored separately.

    \textbf{Protocol.} Five-fold cross-validation grouped by observation is used.
    The score is out-of-fold (OOF) PQ, with TP, FP, FN and IoU pooled over all 1{,}154 records as the evaluator does.
    Post-processing is calibrated on the pooled OOF maps.
    A \emph{nested} score is also reported, with parameters chosen on four folds and applied to the fifth.
    Seeds alone move a five-model set by up to 0.007 and the crop refiner by $\pm0.001$.
    A change is therefore adopted only if it beats a matching control by $>0.002$ after the full pipeline, with the nested score agreeing.


% ==================================== %
%              METHODOLOGY             %
% ==================================== %
\section{Methodology}\label{sec:methodology}


% ==================================== %
%              EVALUATION              %
% ==================================== %
\section{Evaluation}\label{sec:evaluation}




% ==================================== %
%            ACKNOWLEDGMENT            %
% ==================================== %
\section{Acknowledgment}\label{sec:acknowledgment}
    % Do not change
    The organization of the Filament Segmentation Challenge 2026 is partially supported by the U.S. National Science Foundation (NSF) under Grant No. 2209912 and 2433781, directorate for Computer and Information Science and Engineering (CSE), and Office of Advanced Cyberinfrastructure (OAC), and Grant No. 2511630, AST Division Of Astronomical Sciences and MPS Directorate for Mathematical and Physical Sciences.

    The Filament Segmentation Challenge 2026 is sponsored by the U.S. National Science Foundation (NSF) National Solar Observatory (NSO).

    This work utilizes GONG data obtained by the NSO Integrated Synoptic Program, managed by the National Solar Observatory, which is operated by the Association of Universities for Research in Astronomy (AURA), Inc.\ under a cooperative agreement with the National Science Foundation and with contribution from the National Oceanic and Atmospheric Administration. The GONG network of instruments is hosted by the Big Bear Solar Observatory, High Altitude Observatory, Learmonth Solar Observatory, Udaipur Solar Observatory, Instituto de Astrofísica de Canarias, and Cerro Tololo Interamerican Observatory.

\bibliographystyle{ACM-Reference-Format}
\bibliography{main}

\end{document}
